\subsection{Role of Learning}
\label{sec:disc-learning}
At nominal load, learning does not lower the energy bill relative to a
well-designed rule: the price-aware heuristic, which executes the Level-2 plan
per vehicle, is cheaper than the nested controller,% claim: pa_cheaper_than_nested_all_main
\ and neither the learned prices nor the learned dispatch reduce cost relative
to the plan (Table~\ref{tab:decomp}).% claim: decomp_nested_planprice_costlier_than_plan_S3
% claim: decomp_nested_flatprice_costlier_than_plan_S3
\ The learned levels reduce curtailment, which mattered under stress
(S7);% claim: decomp_nested_less_curt_than_plan_S7
\ the lower unmet energy under stress comes from least-laxity-first
allocation, not from learning.% claim: plan_less_unmet_than_pa_S7
% claim: decomp_nested_more_unmet_than_plan_S7
\ Learned dispatch needs the plan as a starting point: learned from scratch,
it is costlier at both penetration levels tested (S3 and
S5).% claim: abl_abl_no_prior_costlier_S3
% claim: abl_abl_no_prior_costlier_S5
\ The nested design is therefore not a way to beat optimization on cost but a
way to keep learned components where they help while planning and physics
handle the rest.

\subsection{Scope of the Voltage Correction}
Level~3 is corrective control on the simulated quasi-static state, not an a
priori guarantee. It restores the floor whenever the reactive headroom of the
plugged-in chargers suffices, and otherwise sheds EV power until the floor is
met or the iteration limit is reached (Section~\ref{sec:method}). Three
consequences follow. First, the inverter apparent-power rating sets the
voltage margin: Section~\ref{sec:res-sens} reports the effect of $S_i$, and S7
the behavior when the reactive capability is halved. Second, we assume, as
in~\cite{leemput2015}, that a charger can supply reactive power at any active
power, including when it does not charge; chargers without this capability
offer only the curtailment fallback. Third, the floor is applied at the
medium-voltage (\SI{12.66}{kV}) buses of the test feeders; drops in service transformers and
secondary circuits are not modeled, and the correction acts at the 60-s
resolution of the simulation, so switching transients and sub-cycle dynamics
are outside its scope, as in any quasi-static study.

\subsection{Assumptions and Limitations}
\emph{Information.} The plan, the learned agents and the heuristic know each
plugged-in vehicle's departure time and remaining need, as in managed-charging
programs where drivers enter them~\cite{lee2019acn}. They do not know future
arrivals. Only the LP-OPF reference sees the whole day in advance.
\emph{Network.} The feeders are balanced, radial and single-phase equivalent;
unbalanced three-phase operation can change the voltage sensitivities on which
the correction step relies.
\emph{Behavior.} The acceptance model~\eqref{eq:fishbein} is stylized. The
voltage layer does not depend on it, but the value of price signals does.
\emph{Fleets.} The residential fleet is synthetic with stated parameters; the
ACN experiments use recorded sessions from workplace sites.
\emph{Data.} Household loads (Pecan Street, United States) and day-ahead prices
(Netherlands) come from different systems; they provide realistic daily shapes,
not a model of one particular utility. Training and test days share the
base-load profile and differ in prices, fleets and noise. Adaptation to
structural market change is not tested beyond the 2019 price regime.
